% Transcription — fonds Alexandre Grothendieck, shelfmark 137, batch 1 (pages 1-20).
% Established from the facsimile archives/batches/137/batch-01.pdf (a copy of
% 137.pdf fetched through the project's relay,
% https://grothendieck-relay.onrender.com/source/137.pdf; PDF page k is
% archivists' page k, no cover sheet), per the skill transcribe-grothendieck.
% This is the first of the folder's five batches (85 pages: 1-20, 21-40,
% 41-60, 61-80, 81-85), transcribed in parallel.
%
% Pass: Opus 5.5 (claude-opus-5-5), 2026-09-26 — first pass, unchecked against the pages by a human.
%
% Pages skipped: 2 and 4, leaves of an English typescript on minimal models
% and de Rham forms (« Appendix H (Hirsch Method) », typed page « H2 »; « The
% Model of a Riemannian Manifold ») carrying nothing in his hand — typed
% versos with nothing of his, skipped by the settled rule; 5, a blank cream
% sheet showing only the show-through of page 4.
%
% Pages summarised: none. Pages 1 and 3 are very faint pencil on thin paper
% (worked through a contrast-stretched red channel); their formulas carry,
% much of the connecting prose does not and is left \ill{}. Pages 6-20 are
% pencil on hole-punched lined sheets (pages 8-12, 14-15, 17-19 on music-staff
% paper) and read nearly end to end.
%
% His pagination: none visible on these leaves. His numbered runs: pp. 1 and
% 3, « Lemme 1 », « Cor », « Lemme 2 », « Théorème »; pp. 6-12, an unnumbered
% « Lemme » (p. 6), then « Prop. » (= Prop. 1, which he cites so), « Cor 1 »,
% « Cor 2 » (p. 8), « Proposition 2 » with its steps 1°) and 2°) (pp. 9-11),
% a cancelled « Lemme », a « Remarque » (p. 11), and a cancelled definition of
% « chain-parfait » modules (p. 12); pp. 13-20, a fresh start: « Déf 1 »,
% « Prop 1 » a)-e) (p. 13), « Th. » (p. 15) with its conditions a)-d) (p. 16),
% « Remarques » 1)-3) (pp. 18-19), and the Hom computations (p. 20). Note that
% p. 6's Lemme cites Cor 1 and 2, which stand on p. 8: the leaves are not in
% his order of writing, and the transcription keeps the archivists' order.
%
% What the batch is about. Semi-simplicial k-modules L. (his « ½simpliciaux »)
% and their Dold-Puppe chain complex L^!, over a commutative ring k.
% Pp. 1, 3: projectivity of a semi-simplicial module over k_* from flasqueness
% plus projective components (Lemme 1), a vanishing of Hom(L. ⊗ Λ^qΨ_*,
% Λ^{q'}Ψ_* ⊗ M) (Lemme 2), and a morphism of complexes from the Koszul-like
% complex Γ^{n-i}Φ_* ⊗ Λ^iΨ_* to the Λ^{i+1}Φ_* with differentials ∧T
% (Théorème). Pp. 6-12: the structure theorem — L. with L_n and π_n(L.)
% projective is a product of Λ^iΦ_{*k} ⊗ M_i and Λ^jΨ_{*k} ⊗ N_j, matching
% the decomposition of L^! into two-term and one-term complexes; the rank
% function rg L_n = Σ μ_i C(n+1,i) + Σ ν_j C(n,j), polynomial growth as the
% criterion for finitely many factors, stability of « spéciaux » modules
% under ⊗ and Λ^k (Prop. 2), and the determination of L. by its rank function
% and the ranks of its π_j (the condition f(-1) = Σ(-1)^j ν_j). Pp. 13-20:
% ∂-pseudo-coherent and ∂-perfect semi-simplicial modules (Déf 1, Prop 1),
% then over a principal ring the classification of strictly ∂-pseudo-coherent
% ones as products of Λ^iΦ., Λ^jΨ_* and new indecomposables Φ(i,ℓ), built from
% k --n--> k and classified by elementary divisors, with the Hilbert function,
% flatness criteria, and the Hom groups between the indecomposables.
%
% Notation fixed once for the batch: L_\bullet for his « L. » (L with a low
% dot), L^{!} for the Dold-Puppe chain complex, \Phi_{*k}, \Psi_{*k} (and
% \Phi_{*}, \Psi_{*}, \Phi_{\bullet} as he drops or changes the subscript),
% \Lambda^{i} for his Λ with the index written above it, \Gamma^{i},
% \binom{n+1}{i} for his binomial columns, « rg » and « rang » as written,
% « t.f. » (type fini), « proj. », « str. », « ps. coh. », « quis » kept.
% His « ½simplicial » is kept with the ½ (semi-simplicial, hand.md §2).
% 𝔅 on p. 15 is set \mathfrak{B}.
%
% Readings to check first: the whole middle of p. 1 and the struck lines of
% p. 3; the marginal note on p. 6 (« on peut éviter cette hypothèse ! »?);
% the inserted words in 2°) on p. 10; « précis » / « strict » over the struck
% « parfait » on p. 8; the circled insertion on p. 15; the ends of the lines
% on p. 14 (« constructibles sur ... ») and the rotated margin there; the
% bracketed interlinear phrases on p. 16; the reading « n' = n/(m,n) » on
% p. 20.

\documentclass[11pt,a4paper]{article}
% Le préambule commun à toutes les transcriptions du fonds Grothendieck.
%
% Il tient en une idée : **l'incertitude se compose, elle ne se résout pas.**
% Une transcription qui lisse un mot douteux a détruit la seule chose pour
% laquelle on la faisait — un lecteur doit pouvoir savoir, à chaque endroit,
% ce qui était sur le feuillet et ce qui a été ajouté. Les six macros qui
% suivent sont donc l'appareil critique, et non de la décoration.
%
% Les mêmes commandes sont reconnues par `scripts/render.mjs`, qui produit la
% vue de lecture du site. Une macro ajoutée ici doit l'être aussi là-bas,
% faute de quoi le rendu échouera bruyamment — ce qui est voulu : un rendu
% silencieusement faux est pire qu'un rendu absent.

\ProvidesPackage{grothendieck}[2026/08/08 Transcriptions du fonds Grothendieck]

\usepackage{amsmath,amssymb,amsthm}
\usepackage{mathrsfs}
\usepackage{stmaryrd}
% Les diagrammes commutatifs sont partout dans ce fonds : c'est la langue
% même de la géométrie algébrique de Grothendieck, et une transcription qui
% les rendrait en prose ne transcrirait rien.
\usepackage{tikz-cd}
% Français par défaut — c'est la langue des feuillets — mais l'anglais est
% chargé aussi : la traduction et le résumé sont des documents anglais, et la
% typographie française (espaces avant les deux-points) y serait une faute.
% Ces documents basculent par \selectlanguage{english} en tête de corps.
\usepackage[english,french]{babel}
\usepackage{fontspec}
\usepackage{geometry}
\geometry{a4paper,margin=25mm,marginparwidth=28mm}
\usepackage{marginnote}
\usepackage{xcolor}
% ulem, not soul. `soul`'s \st and \ul are built on a character-by-character
% scan that XeLaTeX's font machinery breaks: the document fails with an
% undefined control sequence rather than degrading. ulem does the same two jobs
% and is Unicode-engine safe. `normalem` stops it hijacking \emph, which the
% transcriptions use for Grothendieck's own emphasis.
\usepackage[normalem]{ulem}
\usepackage{hyperref}
\hypersetup{colorlinks=true,linkcolor=black,urlcolor=black,pdfborder={0 0 0}}

% --- Métadonnées du lot -----------------------------------------------------
% Reprises telles quelles par le script de rendu, qui les affiche en tête de
% la vue de lecture. Elles fixent ce que le fichier prétend couvrir : sans
% elles, une transcription flotte sans point d'attache dans un fonds de seize
% mille feuillets.

\newcommand{\folder}[1]{\gdef\grothFolder{#1}}
\newcommand{\batch}[1]{\gdef\grothBatch{#1}}
\newcommand{\leaves}[2]{\gdef\grothFirst{#1}\gdef\grothLast{#2}}
\newcommand{\foldertitle}[1]{\gdef\grothFoldertitle{#1}}
\gdef\grothFolder{?}\gdef\grothBatch{?}\gdef\grothFirst{?}\gdef\grothLast{?}\gdef\grothFoldertitle{}

% --- Le résumé --------------------------------------------------------------
%
% Il ouvre la lecture modernisée, sous le titre « Résumé », et il est composé
% en petit. Ce n'est pas un ornement : c'est le seul passage du document qui
% ne soit pas des mathématiques, et un lecteur doit pouvoir le reconnaître
% avant de l'avoir lu — puis le sauter s'il connaît déjà le sujet, ou s'y
% arrêter s'il ne le connaît pas. Le filet qui le suit marque la reprise du
% corps.

\newenvironment{resume}
  {\small\setlength{\parskip}{0.45em}}
  {\par\medskip\hrule\medskip}

% --- Mots-clés ---------------------------------------------------------------
%
% En anglais, en clôture du résumé de la lecture modernisée : le vocabulaire
% moderne sous lequel chercher ce que les feuillets construisent. Le manifeste
% du site (`npm run manifest`) les extrait pour étiqueter la cote entière —
% cette ligne est la source unique des tags, il n'y a pas de fichier à côté.
\newcommand{\keywords}[1]{\par\smallskip\noindent
  {\footnotesize\textsc{Keywords} --- \textit{#1}}\par}

% --- L'appareil critique ----------------------------------------------------

% Le numéro de feuillet, en marge. C'est l'ancre qui permet de régler un
% désaccord sur une formule en regardant *un* feuillet plutôt qu'une chemise
% de six cents. Le site s'en sert aussi pour tourner les pages du fac-similé
% quand on fait défiler la transcription.
\newcommand{\leaf}[1]{%
  \marginnote{\footnotesize\color{gray!70}\textsf{#1}}%
  \label{leaf:#1}\ignorespaces}

% Illisible : on ne devine pas. Un mot inventé qui se lit comme les autres est
% le pire résultat possible.
\newcommand{\ill}{\textcolor{red!60!black}{[\,\ldots\,]}}

% Lecture douteuse : le mot est donné, et signalé comme incertain.
\newcommand{\uncertain}[1]{\uline{#1}}

% Ajout de l'éditeur — une lettre manquante, un mot sous-entendu.
\newcommand{\add}[1]{\textcolor{blue!50!black}{[#1]}}

% Biffé par Grothendieck lui-même. Ce qu'il a rayé fait partie du document :
% une rature dit souvent où la pensée a changé de direction.
\newcommand{\struck}[1]{\sout{#1}}

% Note du transcripteur, sur l'état matériel du feuillet.
\newcommand{\note}[1]{\par\smallskip{\small\color{gray!60!black}\textit{[#1]}}\par\smallskip}

% Note marginale de Grothendieck, distincte du corps du texte.
\newcommand{\marginal}[1]{\par\smallskip\noindent\hspace{1em}%
  {\small\color{gray!60!black}#1}\par\smallskip}

% --- Titre ------------------------------------------------------------------

\AtBeginDocument{%
  \begin{center}
    {\large\bfseries Fonds Alexandre Grothendieck --- cote n\textsuperscript{o}~\grothFolder}\\[2pt]
    {\itshape\grothFoldertitle}\\[4pt]
    {\small feuillets \grothFirst--\grothLast\ (lot \grothBatch)}\\[2pt]
    {\footnotesize\color{gray!60!black}Transcription établie d'après le fac-similé numérisé
     par l'Université de Montpellier.\\
     Les lectures incertaines sont soulignées, les illisibles notées [\,\ldots\,],
     les ajouts entre crochets.}
  \end{center}
  \vspace{1em}}

\endinput


\folder{137}
\batch{1}
\pages{1}{20}
\dating{[vers 1975-1976]}
\watermark{Édition de démonstration}
\foldertitle{Dévissage des complexes ½ simpliciaux ∂ - parfaits et structures multiplicatives… (1975 ou 1976) : notes manuscrites (s.d.)}

\begin{document}

\page{1}
\note{crayon très pâle sur papier mince ; les formules se lisent, une grande
partie de la prose de liaison non. Aucune pagination de sa main sur les
feuillets de ce lot}

$k$ anneau commutatif, $k_{*}$ est l'anneau \uncertain{½simpl.} \ill{}

\textbf{Lemme 1.} Soit $L_{*}$ un $k_{*}$-module ½simplicial. Pour que $L_{*}$
soit un $k_{*}$-Module projectif, il suffit qu'il soit flasque (i.e.
$\pi_{i}(L_{*}) = 0$ $\forall\, i \geqslant 0$) et que les $L_{n}$ soient des
$k$-modules projectifs.
\marginal{\ill{}}\note{quelques signes au crayon dans la marge de gauche,
illisibles}

Le \uncertain{complexe} \ill{} (au sens \ill{} que le complexe de chaînes
(des $k$-modules) associé par Dold-Puppe \ill{} est acyclique et à
composantes projectives. On voit alors que \ill{} est isomorphe \ill{}
\note{la moitié médiane du feuillet (une dizaine de lignes) est presque
effacée ; on y distingue « tel » et une formule $L_{q} \simeq \ldots$, rien
de suivi}

donc on \ill{}

\emph{NB} Il suffit \ill{} $L_{*}$ flasque i.e. $\pi_{i}(L_{*}) = 0$
$\forall\, i \geqslant$ \ill{} \ill{} $\pi_{i}(L_{*}) = 0$ pour $i > 0$ et
$\pi_{0}(L_{*})$ $k$-projectif. (On trouve \ill{} sans \ill{}) En effet,
dans l'argument précédent, il faut de plus tenir compte \ill{} \ill{} qui est
un module projectif. Le plus \ill{}. Or il est \ill{} qu'il \ill{} un
complexe \ill{} de chaînes projectif.

\textbf{Cor.} Les $\Gamma^{i}\Phi_{*k} \otimes_{k} \Lambda\Psi_{*k}$
\add{et $\Lambda^{i}\Phi_{*k} \otimes \Lambda^{j}\Psi_{*k}$} (pour
$i \geqslant 0$ \ill{}) sont des modules ½simpliciaux projectifs.
\note{« et $\Lambda^{i}\Phi_{*k} \otimes \Lambda^{j}\Psi_{*k}$ » est ajouté
au-dessus de la ligne, à l'encre, avec un trait qui le rattache à
l'énoncé}

\page{3}
\textbf{Lemme 2.} Soit $L_{\bullet}$ un $k$-module ½simplicial \ill{}
\struck{flasque}\add{\uncertain{quasi-flasque}}, $q > q' \geqslant 0$, $M$ un
$k$-module. Alors
\[
  \operatorname{Hom}_{\uncertain{\mathrm{Zab}}}\bigl(L_{\bullet} \otimes_{k}
  \Lambda^{q}\Psi_{*k},\ \Lambda^{q'}\Psi_{*k} \otimes_{k} M\bigr) = 0
\]
(on utilise seulement que \ill{} \ill{} en degrés $< q$)
\note{l'indice sous « Hom » est peu lisible ; la parenthèse est reliée par
une accolade au facteur $L_{\bullet} \otimes_{k} \Lambda^{q}\Psi_{*k}$}

\struck{\ill{} le complexe de chaînes associé à $L_{\bullet} \otimes_{k}
\Lambda^{q}\Psi_{*k}$ \ill{} de la forme}\note{trois lignes biffées d'un
trait chacune}

\[
  \cdots \longrightarrow L'_{q+1} \longrightarrow L'_{q} \longrightarrow 0
  \longrightarrow \cdots
\]
\uncertain{Sachant} que $\Lambda^{q'}\Psi_{*k} \otimes_{k} M$ \ill{} $M$
placé en degré $q'$ \add{\uncertain{i.e.} $M[q']$}, \ill{} $q' < q$, il est
évident que $\operatorname{Hom}(L'_{\bullet}, M[q']) = 0$ \ill{} \ill{}
\ill{} degrés \ill{}

\begin{tikzcd}
  L'_{q+1} \arrow[r] \arrow[d] & L'_{q} \arrow[r] \arrow[d] & 0 \\
  0 \arrow[r] & M \arrow[r] & 0
\end{tikzcd}

alors $\pi_{q}(L_{\bullet}) = 0 \Rightarrow L'_{q+1}
\xrightarrow{\ d\ } L'_{q}$ est surjectif, donc $L'_{q} \to M$ est nul,
\ill{} \ill{} \add{\uncertain{de chaînes}} $L'_{\bullet} \to M[q]$, cqfd.

\textbf{Théorème.} Soit $n \geqslant 0$. Il existe un \ill{} de complexes de
\uncertain{cochaînes} dans \ill{}

\begin{tikzcd}[column sep=small, row sep=normal, nodes={font=\scriptsize}]
  0 \arrow[r] & \Gamma^{n}\Phi_{*} \arrow[r] \arrow[d, "u_{0}"] & \Gamma^{n-1}\Phi_{*} \otimes \Psi_{*} \arrow[r] \arrow[d, "u_{1}"] & \Gamma^{n-2}\Phi_{*} \otimes \Lambda^{2}\Psi_{*} \arrow[r] \arrow[d, "u_{2}"] & \cdots \arrow[r] & \Lambda^{n}\Psi_{*} \arrow[r] \arrow[d, "u_{n}"] & 0 \arrow[d] \\
  0 \arrow[r] & \Phi_{*} \arrow[r, "\wedge T"'] & \Lambda^{2}\Phi_{*} \arrow[r, "\wedge T"'] & \Lambda^{3}\Phi_{*} & \cdots \arrow[r, "\wedge T"'] & \Lambda^{n+1}\Phi_{*} \arrow[r, "\wedge T"'] & \Lambda^{n+2}\Phi_{*}
\end{tikzcd}
\note{les exposants des $\Lambda$ de la ligne du bas sont écrits au-dessus
du signe et en partie effacés ; ceux qui sont donnés ici suivent le décalage
d'une unité que portent $u_{0}$ et $u_{1}$. La flèche verticale de droite
n'a pas d'étiquette lisible}

(NB : on écrit $\Phi_{*}$ et $\Psi_{*}$ pour $\Phi_{*k}$, $\Psi_{*k}$, et
$\otimes$ \uncertain{l'omission} de $\otimes_{k}$) tel que $u_{n}$ soit
l'isomorphisme de $\Lambda^{n}\Psi_{*}$ sur
$\operatorname{Ker}\bigl(\Lambda^{n+1}\Phi_{*} \to
\Lambda^{n+1}\Psi_{*}\bigr) \hookrightarrow \Lambda^{n+2}\Phi_{*}$
\uncertain{$= Z(\Lambda^{n+1}\Phi_{*})$}
\note{la fin de la ligne, au bas du feuillet, est très pâle : les exposants
$n+1$, $n+2$ sont lus au-dessus des $\Lambda$ avec doute}

\page{6}
\textbf{Lemme.} Supposons les $L_{n}$ et $\pi_{n}(L_{\bullet})$ projectifs
de t.f. sur $k$, et supposons $k$ connexe (i.e. $\operatorname{Spec}$
connexe) (donc $\forall\, M$ projectif sur $k$, $\operatorname{rg}_{k} M
\in \mathbb{N}$ est défini).
\marginal{\uncertain{on peut éviter cette hypothèse !}}\note{écrit en
oblique dans la marge de gauche, à hauteur de « connexe » ; lecture
incertaine}

Conditions équivalentes :
\begin{itemize}
  \item[a)] $L^{!}$ à degrés bornés.
  \item[b)] \struck{Tous} Les $M_{i}$, $N_{j}$ sont nuls sauf un nombre fini,
  i.e. $L_{\bullet}$ est isomorphe \add{à une} somme finie de modules
  ½simpliciaux $\Lambda^{i}\Phi_{*} \otimes M_{i}$ ($i \geqslant 1$) et
  $\Lambda^{j}\Psi_{*} \otimes N_{j}$ ($j \geqslant 1$) ($M_{i}$ et $N_{j}$
  proj. de t.f.)\note{on lit bien « $j \geqslant 1$ » ici, là où la
  Prop.\ de la page 7 a $j \geqslant 0$}
  \item[b')] (Si $k$ principal, plus généralement si tout module proj. de
  t.f. sur $k$ est libre). $L_{\bullet}$ est isomorphe \add{à une} somme
  finie de modules ½simpliciaux $\Lambda^{i}\Phi_{\bullet}$ ($i \geqslant
  1$), $\Lambda^{j}\Psi_{*}$ ($j \geqslant 0$).
  \item[c)] $\operatorname{rang}_{k} L_{n}$ est une fonction polynomiale en
  $n$ pour $n$ grand.
  \item[d)] $\exists\, N$ \add{et $c > 0$} \struck{tel} que pour $n$ assez
  grand, on ait $\operatorname{rang}_{k} L_{n} \leqslant c\, n^{N}$.
\end{itemize}

[De plus, \struck{p} pour $N$ fixé, la condition que l'on ait
$\operatorname{rang}_{k} L_{n} = O(n^{N})$ \struck{\ill{}} signifie que les
$M_{i}$, $N_{j}$ sont nuls pour $i, j > N$ ;
\struck{\ill{} $N$ \ill{} possible \ill{} \ill{} (pour $M_{N} + N_{N} \neq
0$) $\Leftrightarrow$ b) \ill{} (cor.\ 1 et 2)}]
\note{deux lignes et demie biffées ; « (cor.\ 1 et 2) » est écrit au-dessus
de la fin de la ligne. Le premier mot de la ligne « la condition » est
surchargé}
\marginal{et en \ill{} les \uncertain{rangs} \ill{} $M_{N}$ \ill{}
$= \mu_{N} + \nu_{N}$}\note{note de la marge de gauche, écrite en oblique et
reliée par un grand arc au crochet}

\emph{Dém.} \struck{Équivalence de a) b)} (ci-dessus), b)
$\Leftrightarrow$ b') trivial, \struck{a) $\Rightarrow$} c) $\Rightarrow$
d) trivial, a) $\Rightarrow$ c) et d) $\Rightarrow$ a) immédiat par la
formule explicite.

\page{7}
\textbf{Prop.} Soit $L_{\bullet}$ un $k$-module ½simplicial. Les conditions
suivantes sont équivalentes :
\begin{itemize}
  \item[a)] Les $L_{n}$ et $\pi_{n}(L_{\bullet})$ sont des $k$-modules
  projectifs.
  \item[a')] Les composantes \struck{de} $(L^{!})_{n}$ de $L^{!}$ (le
  complexe de chaînes de Dold-Puppe associé à $L_{\bullet}$) et les
  $H_{n}(L^{!})$ sont des $k$-mod.\ projectifs.\note{« a') » surcharge un
  « b) »}
  \item[b)] $L_{\bullet}$ est isomorphe à un \struck{somme}
  \add{produit} \struck{direct de modules} de la forme
  \[
    \Bigl(\prod_{i \geqslant 1} \Lambda^{i}\Phi_{*k} \otimes_{k}
    M_{i}\Bigr) \times \prod_{j \geqslant 0} \bigl(\Lambda^{j}\Psi_{*k}\bigr)
    \otimes_{k} N_{j} .
  \]
  \item[b')] $L^{!}$ est isomorphe à un produit (de complexes de
  \uncertain{chaînes}) de la forme
  \[
    \prod_{i \geqslant 1} \bigl(0 \to M_{i} \to M_{i} \to 0 \to
    \cdots\bigr) \times \prod_{j \geqslant 0} \bigl(0 \to N_{j} \to 0 \to
    \cdots\bigr)
  \]
  (le premier facteur en degrés $i$ et $i-1$, le second en degré
  $j$).\note{« b') » surcharge un « a) » ; les indications de degré sont
  portées sous les facteurs, avec des accolades}
\end{itemize}

De plus, les $M_{i}$ et $N_{j}$ de b) (resp.\ b')) sont déterminés à isom.\
près par $L_{\bullet}$, et ce sont les mêmes (à isom.\ près) dans b) et b').

\emph{Dém.} a') $\Leftrightarrow$ b') est trivial et bien connu, ainsi que
l'unicité à isom.\ \add{can.} près des $M_{i}$, $N_{j}$ dans b') (NB on a
$N_{j} \simeq H_{j}(L_{\bullet})$ et $M_{i} \simeq L_{i}/Z_{i}$ ($\simeq
B_{i-1}$)). On a a) $\Leftrightarrow$ a') et b) $\Rightarrow$ b')
\add{\uncertain{par le}} \ill{} de Dold-Puppe. D'où l'équivalence des 4
conditions, et la dernière assertion est connue.

\page{8}
\textbf{Cor 1.} \struck{Dans ces conditions} On obtient \ill{} des
conditions équivalentes en exigeant \add{de plus} dans a) que les $L_{n}$
\struck{$\pi_{n}(L_{\bullet})$} soient de type fini, dans a') que les
$L^{!}_{n}$ \struck{$H_{n}(L^{!})$} sont de type fini, dans b) et b') que
les $M_{i}$ et $N_{j}$ soient de type fini.

On a alors, dans les conditions de Prop 1,
\[
  L_{n} = \prod_{1 \leqslant i \leqslant n+1} \Lambda^{i}\Phi_{n} \otimes_{k}
  M_{i} \times \prod_{0 \leqslant j \leqslant n} \Lambda^{j}\Psi_{n}
  \otimes N_{j}
\]
\struck{dans un \ill{}}

\textbf{Cor 2.} Sous les conditions équivalentes de \ill{}
\struck{\ill{}}\note{un bloc de deux lignes est en partie biffé et
surchargé ; on n'en lit que la fin} \ill{} degrés bornés, et finis), qu'il
n'y ait qu'un nombre fini d'indices $i$ et $j$ tels que $M_{i} \neq 0$,
$N_{j} \neq 0$.

\struck{Sous les conditions} On dit que $L_{\bullet}$ est
\struck{parfait} \add{\uncertain{spécial}} si $L^{!}$ est à degrés bornés et
composantes \ill{} à $H$ projectifs de t.f.\ — i.e.\ si $L_{\bullet}$ est
isomorphe à somme finie de modules de la forme $\Lambda^{i}\Phi_{*} \otimes_{k}
M_{i}$ ($i \geqslant 1$) et $\Lambda^{j}\Psi_{*} \otimes M_{j}$ ($j
\geqslant 0$), avec les $M_{i}$ et $N_{j}$ proj.\ de t.f., strictement
spécial si de plus les $M_{i}$, $N_{j}$ sont libres. Je reviens au cas de
\note{le mot écrit au-dessus de « parfait » biffé est peu lisible ; la
page suivante emploie « spéciaux » et « strict.\ spéciaux », ce qui fait
lire « spécial »}

\page{9}
dire que $L_{\bullet} \simeq L^{\circ}_{\bullet} \otimes_{\mathbb{Z}} k$, où
$L^{\circ}_{\bullet}$ est un complexe de $\mathbb{Z}$-modules ½simpliciaux
spécial.

\textbf{Proposition 2.} Les $k$-modules \add{½simpliciaux} spéciaux
(resp.\ strict.\ spéciaux) sont stables par produits tensoriels
\add{et par les op.\ tensorielles $\Lambda^{i}$, $\Gamma^{i}$,
\uncertain{$\mathrm{Sym}^{i}$}}.\marginal{\ill{}}

Le cas strictement spécial se ramène au cas spécial, en se plaçant sur
$\mathbb{Z}$ \ill{}

1°) \emph{Produits tensoriels.} On est réduit à prouver que les modules
$\Lambda^{i}\Phi_{*} \otimes \Lambda^{j}\Phi_{*}$ ($i, j \geqslant 1$),
$\Lambda^{i}\Psi \otimes \Lambda^{j}\Psi$ ($i, j \geqslant 0$),
$\Lambda^{i}\Phi \otimes \Lambda^{j}\Psi$ ($i \geqslant 1$, $j \geqslant 0$)
sont spéciaux. Or ils sont à composantes libres de t.f., donc par
Künneth-Eilenberg-Zilber leurs $\pi_{i}$ sont itou, il reste à montrer qu'il
n'y a qu'un nombre fini de facteurs $M_{i}$, $N_{j} \neq 0$.
\struck{On a} Cela résulte \struck{\ill{}} aussitôt du

\struck{\textbf{Lemme.} \struck{Soit} Les conditions du Cor 2 et Prop 1
sont aussi équivalentes à la \add{suivante} (en supposant \add{que $k$ est
connexe}, i.e.\ $\operatorname{Spec} k$ connexe) : Soit $\mu_{i} =
\operatorname{rg} M_{i}$ et $\nu_{j} = \operatorname{rg} N_{j}$ \ill{}
défini) a) $\operatorname{rg}_{k} L_{n}$ est un polynôme en $n$ pour $n$
grand b) $\exists$}\note{le Lemme est barré de traits en diagonale}

\page{10}
\[
  \rho(n) \overset{\mathrm{df}}{=} \operatorname{rg}_{k} L_{n} =
  \sum_{1 \leqslant i \leqslant n+1} \mu_{i} \binom{n+1}{i} +
  \sum_{0 \leqslant j \leqslant n} \nu_{j} \binom{n}{j}
\]
\note{sous les deux termes, des flèches renvoient à « pol.\ de degré $i$ en
$n$ » et « pol.\ de degré $j$ en $n$ »}
d'où $\mu_{i} = \operatorname{rg}_{k} M_{i}$, $\nu_{j} =
\operatorname{rg}_{k} N_{j}$, donc si \struck{$\mu_{i}$ ($\nu_{j}$)} un
$\mu_{i}$ ou un $\nu_{j}$ est $\neq 0$, cela montre que $\rho(n)$ pour $n$
grand est minoré par un polynôme de degré $i$ resp.\ $j$.

2°) \emph{Cas de $\Lambda^{k}$.} Utilisons la formule du $\Lambda^{k}$ d'une
somme directe, et 1°), on est ramené à prouver que le $\Lambda^{k}$ d'un
$k$-module ½simplicial $L_{\bullet}$ de la forme $\Lambda^{i}\Phi \otimes M$
ou $\Lambda^{j}\Psi \otimes N$ \ill{}. Comme $M$, $N$ sont facteurs directs
d'un libre \add{\uncertain{$P$}}, \add{$L_{\bullet}$ est facteur direct de
$L'_{\bullet} = \Lambda^{i}\Phi \otimes P$ ou $\Lambda^{j}\Psi \otimes P$,}
\struck{\ill{}} \add{cela, grâce à Cor.\ 1 et 2 \add{du Cor 2
\struck{\ill{}} est stable par facteurs directs}, et \ill{}} $\Lambda^{k}$
d'une somme directe \ill{} montre que $\Lambda^{k}(L_{\bullet})$ est facteur
direct de $\Lambda^{k}L'_{\bullet}$, on est ramené à prouver
$\Lambda^{k}L'_{\bullet}$ spécial, donc \ill{} on peut supposer $M$, $N$
libres — donc (par la formule de la somme) $M = k$, $N = k$. Dans ce cas on
\uncertain{est} ramené à prouver
\note{les ajouts interlinéaires de ce paragraphe se chevauchent et sont
reliés par des boucles ; leur ordre ici est une reconstitution}

\page{11}
$\Lambda^{k}(\Lambda^{i}\Phi_{*k})$ et $\Lambda^{k}(\Lambda^{j}\Psi_{*k})$
sont spéciaux (bien sûr, ils seront même très spéciaux !). Leurs
composantes sont évidemment libres. De plus, comme $\Lambda^{i}\Phi_{*k}$
est homotope à $0$, de même $\Lambda^{k}(\Lambda^{i}\Phi_{*k})$, donc
\ill{} les $\pi_{n}$ sont des modules \ill{} projectifs de t.f.]

\textbf{Remarque.} Soit $L_{\bullet}$ un $k$-module ½simplicial spécial,
alors $L_{\bullet}$ est connu à isom.\ près, i.e.\ les $\mu_{i} =
\operatorname{rg} M_{i}$ ($i \geqslant 1$) et $\nu_{j} = \operatorname{rg}
N_{j}$ \add{($j \geqslant 0$)} \struck{sont} connus, quand on connaît la
\struck{\ill{}} fonction $n \mapsto \operatorname{rg}_{k} L_{n} =
f_{L_{\bullet}}(n)$ \struck{\ill{}} et les $\operatorname{rg}_{k}
\pi_{j}(L_{\bullet})$ ($= \operatorname{rg}_{k} N_{j} = \nu_{j}$). En
effet, on aura
\[
  \sum_{i \geqslant 1} \mu_{i} \binom{n+1}{i} = f_{L_{\bullet}}(n) -
  \sum_{j \geqslant 0} \nu_{j} \binom{n}{j} \qquad (n \geqslant -1)
\]
\note{au-dessus de $\binom{n+1}{i}$ est écrit « $= P_{i}(n+1)$ »}
\struck{On en sait} i.e.
\[
  f_{L_{\bullet}}(n-1) - \sum_{j \geqslant 0} \nu_{j} P_{j}(n-1) =
  \sum_{i \geqslant 1} \mu_{i} P_{i}(n)
\]
et on sait que les $P_{i}(n)$ \ill{} sont \uncertain{linéairement}
indépendants — et forment une base du $\mathbb{Z}$-module des polynômes
\ill{} à valeurs entières sur les entiers. Donc la connaissance de
$\sum_{i \geqslant 1} \mu_{i} P_{i}$ pour $n$ grand implique

\page{12}
la connaissance des $\mu_{i}$. Si on se donne des $\nu_{j} \geqslant 0$
\add{entiers} \struck{(${}= \operatorname{rg} \pi_{j}(L_{\bullet})$)}
($j \geqslant 0$) et une fonction \add{polynomiale} $f(n)$ à valeurs
entières \struck{définie pour $n$ grand} \add{sur les entiers \ill{}},
condition nécess.\ et suff.\ pour qu'il existe $L_{\bullet}$ très spécial
avec $\nu_{j} = \operatorname{rg} \pi_{j}(L_{\bullet})$ ($j \geqslant 0$),
et $f(n) = f_{L_{\bullet}}(n)$ pour $n$ grand ? Notons que
$g = f - \sum_{j \geqslant 0} \nu_{j} P_{j}$ est encore une fonction
\struck{entière} polynomiale (à valeurs entières), donc de la forme
$\sum_{i \geqslant 0} \mu_{i} \binom{n+1}{i}$ où $(\mu_{i}) \in
\mathbb{Z}^{(\mathbb{N})}$ est uniquement déterminé. Alors la condition
cherchée est $\mu_{0} = 0$, $\mu_{i} \geqslant 0$ si $i \geqslant 1$.
Notons que \struck{$\mu_{0}$} $\mu_{0} = g(-1) = f(-1) - \sum_{j \geqslant 0}
\nu_{j} \frac{(-1)(-2)\cdots(-j)}{j!}$, donc on a
\[
  (\mu_{0} = 0) \iff \Bigl(f(-1) = \sum (-1)^{j}\nu_{j}\Bigr)
  \quad \bigl[= \chi((\nu_{j}))\bigr].
\]

\struck{\textbf{Déf.}} \struck{Complexes} \struck{\ill{}}
\struck{Modules ½simpliciaux} \struck{$L_{\bullet}$ \ill{}} \add{str.}
\struck{« chain-parfait »} \add{(strictly chain-perfect)} \struck{$=$ tels
que} \struck{$L^{!}$ soit} \struck{un complexe} \struck{de chaînes}
\struck{de $k$-modules} \struck{str.} \struck{parfait} \struck{(str.\
parfait) $=$}\note{ce début de
définition, qui s'interrompt, est barré de quatre traits en diagonale}

\page{13}
\note{ici commence, sur un papier réglé ordinaire, une nouvelle rédaction,
qui reprend la numérotation : Déf 1, Prop 1}
$k$ anneau (comm.\ unit.)

\textbf{Déf 1.} Un $k$-module ½simplicial $L_{\bullet}$ est dit
$\partial$-pseudo-cohérent (resp.\ strictement $\partial$-pseudocohérent,
resp.\ $\partial$-parfait, resp.\ strictement $\partial$-parfait) si
$L^{!}$ (le complexe de chaînes associé par Dold-Puppe) est
pseudo-cohérent (resp.\ str.\ ps.\ cohérent, resp.\ parfait, resp.\
strictement parfait). Donc on a par déf.\ a) ci-dessous.

\textbf{Prop 1.}
\begin{itemize}
  \item[a)] $L_{\bullet}$ est str.\ $\partial$-pseudocohérent (resp.\ str.\
  $\partial$-parfait) si et seulement si les $L^{!}_{n}$ sont projectifs de
  t.f.\ (resp.\ et nuls sauf un nb fini).
  \item[b)] $L_{\bullet}$ est str.\ $\partial$-ps.\ coh.\ ssi les $L_{n}$
  sont projectifs de t.f.
  \item[c)] $L_{\bullet}$ est str.\ $\partial$-parfait ssi $L_{\bullet}$ est
  extension successive \add{finie} de Modules de la forme
  $\Lambda^{i}\Psi_{\bullet} \otimes_{\mathbb{Z}} M_{i}$ ($i \geqslant 0$),
  où les $M_{i}$ sont $k$-modules pr.\ de t.f.\ (l'ordre des facteurs dans
  une telle filtration pouvant d'ailleurs être choisi par ordre de $i$
  croissant), ou encore \struck{ssi} (si on veut $k$ noeth.\ ou simplement
  \ill{})
  \marginal{ssi les $L_{n}$ sont proj.\ de t.f.\ et
  $f_{L_{\bullet}}(n) = \operatorname{rg}_{k} L_{n}$ est une fonction
  polynomiale pour $n$ grand (i.e.\ \ill{} telles fonctions)}\note{la
  marge de gauche porte ce complément, écrit en oblique vers le bas et relié
  à « ou encore »}
  \item[d)] Si $k$ est noeth., $L_{\bullet}$ est $\partial$-pseudo-cohérent
  (ssi les $L^{!}_{n}$ sont de type fini sur $k$, ou encore) ssi les $L_{n}$
  le sont.
  \item[e)] Si $k$ est noeth.\ régulier, $L_{\bullet}$ est
  $\partial$-parfait ssi les $L_{n}$ sont de t.f.\ et \struck{\ill{}}
  \ill{} évident, la fonction
\end{itemize}

\page{14}
$n \mapsto \operatorname{rg}_{K} L_{n} \otimes_{k} K$ est polynomiale
\struck{pour} en $n$ grand \add{(supposé non nul, i.e.\ $L_{\bullet} \neq
0$)}, de degré du polynôme (est alors le plus grand entier $N$ tel que
$L^{!}_{N} \neq 0$), et les coefficients du polynôme exprimé comme
combinaison linéaire des $P_{i}(n) = \binom{n}{i}$ ($i \geqslant 0$) sont
les $\operatorname{rg}_{K} L^{!}_{i} \otimes_{k} K$, les coeffs dudit
polynôme (\ill{} avec l'exception habituelle \struck{\ill{}} \ill{}) sont
des fonctions constructibles sur le \ill{} \uncertain{($K$ corps résiduel
en un point de $\operatorname{Spec} k$)}.

\emph{NB} Si dans e) on suppose $k$ noeth.\ mais non régulier, la condition
\add{énoncée} de $\partial$-perfection est néc., mais pas suffisante — il
faut une condition néc.\ et suff.\ plus subtile, savoir l'existence d'un
$L'_{\bullet} \to L_{\bullet}$, avec $L'_{\bullet}$ str.\ $\partial$-parfait
(cf.\ c)) et qui soit un quis, i.e.\ induise un isom.\ sur les $\pi_{i}$.
(De plus, alors $L'_{\bullet} \to L_{\bullet}$ épi.)

\marginal{Stabilité des notions par extensions, \uncertain{produits}
\ill{}}\note{écrit verticalement le long de la marge de gauche, au bas de la
page}

On veut une théorie \ill{} des $L_{\bullet}$ (str.\ $\partial$-ps.\ coh.)
\struck{\ill{}} — cela n'est possible bien sûr que
\struck{\ill{} stabilité de \ill{} $\partial$-ps.\ coh.}

\page{15}
\struck{\ill{}} si on dispose d'une théorie analogue pour les complexes de
chaînes \struck{\ill{}} ps.\ coh., ce qui n'est guère le cas que si
\struck{\ill{}} $k$ principal (\uncertain{ou} Dedekind à la rigueur...).
Supposons donc $k$ principal \struck{\ill{}} (le cas $k$ un corps est déjà
traité avec les \ill{} des $L_{\bullet}$ très spéciaux \struck{\ill{}} et
str.\ spéciaux), soit $\operatorname{Max}(k)$ l'ens.\ des idéaux
\struck{premiers} maximaux de $k$, $\mathfrak{B} =
\mathbb{N}^{(\operatorname{Max} k)} - (0)$
\add{\ill{} d'un projectif de t.f.\ est \ill{}}\note{insertion entourée,
au-dessus de la ligne, peu lisible}

\textbf{Th.} Tout $L_{\bullet}$ \add{str.} $\partial$-ps.\ coh.\ est
isomorphe à un produit de modules de la forme suivante
\begin{itemize}
  \item[a)] $(\Lambda^{i}\Phi_{\bullet})^{\mu_{i}}$ ($i \geqslant 1$)
  \add{$\mu_{i} \in \mathbb{N}$} \struck{$M_{i}$ $k$-module libre de t.f.}
  \item[b)] $(\Lambda^{j}\Psi_{*})^{\nu_{j}}$ ($j \geqslant 0$) ($\nu_{j} \in
  \mathbb{N}$)
  \item[c)] \struck{\ill{}} $\Phi(i, \ell)^{\rho_{i,\ell}}$ ($i \geqslant 1$)
  $\rho_{i,\ell} \in \mathbb{N}$ \struck{où $\Phi(i,\ell)$ \ill{}} $\ell$ un
  \struck{idéal} diviseur de $k$ (cf.\ plus bas la définition de
  $\Phi(i,\ell)_{*}$),
\end{itemize}
de façon précise
\[
  L_{\bullet} \simeq \prod_{i \geqslant 1} M_{i} \otimes_{k}
  \Lambda^{i}\Phi_{\bullet} \times \prod_{j \geqslant 0} N_{j} \otimes_{k}
  \Lambda^{j}\Psi_{*} \times \prod_{\substack{i \geqslant 1\\ \ell \in
  \mathfrak{B}}} \Phi_{i,\ell}^{\rho_{i,\ell}}
\]
où pour tout \struck{\ill{}} $i \geqslant 1$, \uncertain{les}
$\rho_{i,\ell}$ sauf un nb fini sont nuls. (NB de tels produits sont bien
str.\ $\partial$-ps.\ coh.) On peut supposer de plus que pour tout $i$
donné, l'ens.\ des $\ell \in \mathfrak{B}$ \struck{\ill{}} tels que
$\rho_{i,\ell} \neq 0$ soit \uncertain{totalement} ordonné, moyennant

\page{16}
ces conditions, les entiers $\mu_{i}$ ($i \geqslant 1$), $\nu_{j}$ ($j
\geqslant 0$), $\rho_{i,\ell}$ ($i \geqslant 1$, $\ell \in \mathfrak{B}$) sont
unique\supplied{ment} déterminés par les conditions suivantes
\begin{itemize}
  \item[a)] Pour tout $i \geqslant 1$, les $\rho_{i,\ell} \neq 0$ sont les
  \add{multiplicités des} diviseurs élémentaires de \struck{Tors}
  $\pi_{i-1}(L_{\bullet})$ ($= H_{i-1}(L^{!})$)
  \item[b)] \struck{\ill{}} $\nu_{i} = \operatorname{rg} \pi_{i}(L_{\bullet})$
  ($= \operatorname{rg} H_{i}(L^{!})$)
  \item[c)] On a \struck{rang} $\mu_{i} = \operatorname{rg}
  B_{i-1}(L^{!})$ \struck{\ill{}} $- \rho_{i}$ (où $\rho_{i} =
  \sum_{\ell} \rho_{i,\ell}$)
  \item[d)] On a de plus
  \[
    \begin{cases}
      \operatorname{rg} Z_{i}(L^{!}) = \mu_{i+1} + \rho_{i+1} + \nu_{i} \\
      \operatorname{rg} L^{!}_{i} = \mu_{i+1} + \rho_{i+1} + \rho_{i} +
      \nu_{i}
    \end{cases}
  \]
  \note{la seconde ligne devrait contenir aussi $\mu_{i}$, par c) ; le
  manuscrit ne l'écrit pas, ou il est effacé sous une surcharge}
\end{itemize}
\add{Enfin $L_{\bullet}$ est (str.) $\partial$-parfait ssi le nb de
facteurs non nuls qui interviennent est fini.}\note{ajouté entre les
lignes}

\emph{Dém.} \struck{\ill{}} On se ramène en \uncertain{vertu} de
Dold-Puppe, à une question \add{tout à fait} évidente sur les complexes de
chaînes. \ill{} \struck{\ill{}} Un sous-module d'un $k$-module projectif de
t.f.\ \ill{} est itou, on conclut que tout tel complexe, s'il est
\struck{pseudo} str.\ ps.\ coh.\ \ill{} à comp.\ proj.\ de t.f., est
\struck{\ill{}} \ill{} produit de complexes (définis à isom.\ \add{non}
unique près, \struck{par le choix des} \ill{}) de la forme
\[
  L^{!} \simeq \prod_{i} K_{i}, \qquad
  K_{i} : \quad \cdots \to 0 \to P_{i} \hookrightarrow Q_{i} \to 0 \to
  \cdots \quad \text{(en degrés $i$ et $i-1$)}
\]
avec $P_{i} = B_{i-1} \hookrightarrow Z_{i-1} \simeq Q_{i}$.
\note{les indications « degré $i$ », « degré $i-1$ » sont portées
au-dessus des deux termes}

\page{17}
On désigne par $\overline{P}_{i} \subset Q_{i}$ \struck{le \ill{}}
\ill{} le sous-gpe \ill{} de torsion de $Q_{i}/B_{i} \simeq H_{i-1}$ ;
$N_{i-1} = Q_{i}/\overline{P}_{i}$ est de t.f.\ sans torsion, donc
\struck{$H_{i}$/torsion} projectif de t.f., donc $Q_{i} \simeq
\overline{P}_{i} \oplus N_{i-1}$ et le complexe $K_{i}$ est isom.\ à la
somme directe des complexes
\[
  \overline{K}_{i} : \ \cdots \to 0 \to P_{i} \hookrightarrow
  \overline{P}_{i} \to 0 \to 0, \qquad
  \cdots \to 0 \to N_{i-1} \to 0
\]
\note{les degrés $i$, $i-1$ sont notés au-dessus des termes, comme à la page
précédente}
Enfin, $P_{i} \hookrightarrow \overline{P}_{i}$ est une inclusion de
modules projectifs de t.f.\ tels que $\overline{P}_{i}/P_{i}$ ($\simeq
\operatorname{Tors} H_{i-1}$) soit de torsion, d'où justiciable de la
théorie des diviseurs élémentaires \add{si $k$ principal} :
\struck{\ill{}} le complexe est une somme directe \struck{$K_{i} \simeq$} de
complexes de la forme
\[
  \cdots \to k \xrightarrow{\ m_{i,\alpha}\ } k \to 0 \qquad
  (1 \leqslant \alpha \leqslant N_{i},\ N_{i} = \operatorname{rg} P_{i} =
  \operatorname{rg} \overline{P}_{i})
\]
avec les $m_{i,\alpha} \in k - (0)$ ($m_{i,\alpha}$ \struck{\ill{}}
\add{définis à} une unité \struck{\ill{}} près) ; le nombre des
$m_{i,\alpha}$ qui sont des unités est $\mu_{i}$ (cela leur correspond
$(\Lambda^{i}\Phi)^{\mu_{i}}$), et si on prend les $m_{i,\alpha}$ tels qu'ils
soient croissants pour la divisibilité, on trouve une suite unique ; les
multiplicités d'occurrences des diviseurs \ill{}

\page{18}
dans la suite des $\operatorname{div}(m_{\alpha})$ est $\rho_{i,\ell}$, et
$\Phi(i,\ell)_{*}$ est le \struck{complexe} \add{module ½simplicial}
associé à
\[
  \cdots \to 0 \to k \xrightarrow{\ n\ } k \to 0 \to \cdots \qquad
  \text{(en degrés $i$ et $i-1$)}
\]
où $\operatorname{div}(n) = \ell$. On a donc une extension
\[
  0 \to \Lambda^{i-1}\Psi \to \Phi(i,n) \to \Lambda^{i}\Psi \to 0
\]
où l'homom.\ assoc.\ à cette extension $k \to k$ (cf.\ \ill{}) est
$n \cdot \mathrm{id}$, où $\operatorname{div}(n) = \ell$, donc
\struck{\ill{}}
\[
  \pi_{j}(\Phi(i,n)) =
  \begin{cases}
    0 & \text{si } j \neq i-1 \\
    k/\operatorname{div} n & \text{si } j = i-1 .
  \end{cases}
\]

\textbf{Remarques.} 1) La fonction de Hilbert de $\Phi(i,n)_{*}$ est donc
la même que celle de $\Lambda^{i}\Phi_{\bullet}$, savoir
$\binom{n+1}{i}$. Donc si $L_{\bullet}$ est \add{str.}
$\partial$-parfait, il \struck{\ill{}} est donné par la \ill{} de la page
\ill{} :
\[
  \operatorname{rg} L_{n} = \sum_{i \leqslant n+1} (\mu_{i} + \rho_{i})
  \binom{n+1}{i} + \sum_{0 \leqslant j \leqslant n} \nu_{j} \binom{n}{j} .
\]
2) \struck{\ill{}} Donc $L_{\bullet}$ est connu à isom.\ \add{non unique}
près quand on connaît a) $\operatorname{rg} \pi_{i}(L_{\bullet}) =
\nu_{i}$ et le nombre $\rho_{i+1}$ de ses diviseurs de torsion, pour tout
$i \geqslant 0$, et b) la fonction $f_{L_{\bullet}}$ pour $n$ grand (qui
s'écrit \struck{alors} \add{pour $n$ grand} comme le polynôme
$\sum_{i \geqslant 1} (\mu_{i} + \rho_{i}) \binom{n}{i} + \sum_{j \geqslant
0} \nu_{j} \binom{n}{j}$\note{ici $\binom{n}{i}$, là où la formule de 1)
porte $\binom{n+1}{i}$ ; la parenthèse n'est pas refermée}

\page{19}
2) Pour que $L_{\bullet}$ soit \struck{projectif} \add{\uncertain{plat}
(projectif)}, il f.\ et s.\ qu'il n'y ait que des facteurs
$(\Lambda^{i}\Phi_{*})^{\mu_{i}}$ (i.e.\ \struck{\ill{}}
$(\Lambda^{0}\Psi_{*})^{\nu_{0}} = k^{\nu_{0}}$).
\add{Pour que $L_{\bullet} \otimes_{k} K$ ($K$ corps des fractions) soit
\ill{} il f.\ et suffit qu'il n'y ait que des
$(\Lambda^{i}\Phi_{\bullet})^{\mu_{i}}$, $\Phi_{*}(i,n)^{\rho_{i,\ell}}$
($i \geqslant 1$) à l'exclusion des $(\Lambda^{j}\Psi_{*})^{\nu_{j}}$ ($j
\geqslant 0$).}\note{ajout entouré, entre les lignes. Il numérote 2) deux
fois : à la fin de la page 18 et ici}

\struck{Pour} 3) Pour déterminer \struck{\ill{}} la catégorie additive des
$k$-modules ½simpliciaux str.\ $\partial$ ps.\ coh.\ (resp.\ str.\
$\partial$ parfaits), il suffit en principe de déterminer les homom.\ entre
les objets indécomposables $\Lambda^{i}\Phi$, $\Phi(i,n)$ \add{$n$ pas
unité, défini à \ill{} près, défini par $\operatorname{div}(n) = \ell$},
$\Lambda^{j}\Psi_{*}$ ($i \geqslant 1$, $j \geqslant 0$) — qu'on peut
réduire à deux types si on admet aussi le cas $n$ unité : $\Phi(i,n)$,
$\Lambda^{j}\Psi_{*}$ ($i \geqslant 1$, $j \geqslant 0$). On trouve, en
repassant aux complexes de chaînes associés

\page{20}
On \struck{a} \ill{}
\[
  \operatorname{Hom}_{k}(\Phi(i,n), L_{\bullet}) \simeq
  d_{L_{\bullet},i}^{-1}(n L_{i-1}) \subset L_{i}
\]

\begin{tikzcd}
  \cdots \arrow[r] & k \arrow[r, "n"] \arrow[d, "u_{i}"] & k \arrow[r] \arrow[d, "u_{i-1}"] & \cdots \\
  \cdots \arrow[r] & L_{i} \arrow[r, "d_{i}"] & L_{i-1} \arrow[r] & \cdots
\end{tikzcd}

($\xi_{i} \mapsto n\,\xi_{i-1}$), donc
\struck{$\operatorname{Hom}(\Phi(i,n), \Phi_{j}(m)) = \{0$ si $j \neq i,
i+1$}
\[
  \operatorname{Hom}(\Phi(i,n), \Phi(j,m)_{*}) \simeq
  \begin{cases}
    0 & \text{si } j \neq i,\ i+1 \\
    \{\xi \in k \mid m\xi \in nk\} = n'k \simeq k & \text{si } j = i \\
    k & \text{si } j = i+1
  \end{cases}
\]
\note{« (où $n' = n/(m,n)$) » est ajouté à droite du cas $j = i$ ; lecture
incertaine}
\add{engendré par le composé $\Phi(i,n) \xrightarrow{\mathrm{can}}
\Lambda^{i}\Psi_{*} \xrightarrow{\mathrm{can}} \Phi(i+1,m)_{*}$} [p.\ ex.\
on trouve $\operatorname{Hom}_{k}(\Phi(i,n)_{*}, \Phi(i,n)_{*}) \simeq k
\cdot \mathrm{id}_{\Phi(i,n)} \simeq k$ (isom.\ \struck{\ill{}} de
$k$-algèbres) \ill{} $\forall\, i \geqslant 0$, $n \in k - \lbrace 0
\rbrace$. \struck{$\operatorname{Hom}(\Lambda^{i}\Phi, \Lambda^{i+1}\Phi)
\ldots$} En particulier
$\operatorname{Hom}_{k}(\Lambda^{i}\Phi_{*}, \Lambda^{i}\Phi_{*}) \simeq
k\, \mathrm{id} \simeq k$, $\operatorname{Hom}_{k}(\Lambda^{i}\Phi_{*},
\Lambda^{i+1}\Phi_{*}) \simeq k \cdot \lambda \simeq k$, où $\lambda =
\uncertain{T_{\wedge}} : \Lambda^{i}\Phi_{*} \to \Lambda^{i+1}\Phi_{*}$]
\[
  \operatorname{Hom}_{k}(\Phi(i,n), \Lambda^{j}\Psi_{*}) \simeq
  \begin{cases}
    0 & \text{si } j \neq i \\
    k & \text{si } j = i
  \end{cases}
\]
[en particulier
\[
  \operatorname{Hom}_{k}(\Lambda^{i}\Phi_{*}, \Lambda^{j}\Psi_{*}) \simeq
  \begin{cases}
    0 & \text{si } j \neq i \\
    k & \text{si } j = i
  \end{cases}
\]
et le générateur dans le cas $j = i$ est l'homom.\ can.\
$\Lambda^{i}\Phi_{*} \xrightarrow{\mathrm{can}} \Lambda^{i}\Psi_{*}$]
\[
  \operatorname{Hom}_{k}(\Lambda^{i}\Psi_{*}, L_{\bullet}) \simeq
  Z_{i}(L_{*}) \simeq k^{\mu_{i+1} + \rho_{i+1} + \nu_{i}}
\]
en particulier

\end{document}
